
% To be included
\chapter*{现代化、技术革命与控制论}

\footnote{本文为《工程控制论》(修订版)(1980年)的序言。}第一版《工程控制论》原是用英文写的，出版于1954年\footnote{Tsien, H. S., \textit{Engineering Cybernetics}. McGraw-Hill Book Company, (1954).}，俄文版是1956年\footnote{\rufont{Цянь-Сюэ-Сянь. Техническая Кибернетика, перевод с английского М. З. Литвина-Селого, под редакцией А. Фельбаума, Изв. Иностранной Литературы, (1956).}}，德文版是1957年\footnote{Tsien, H. S., \textit{Technische Kybernetik}, übersetzt von Dr. H. Kaltenecker, Berliner Union, (1957).}，中文版是1958年\footnote{钱学森：《工程控制论》，戴汝为等译自英文版，科学出版社，1958年。}。现在回顾那个时代，恍同隔世！在这20多年中，我们的国家和整个世界都经历了天翻地覆的变化。我国人民经受了这一伟大时代的革命锻炼，正走上新的长征，为实现四个现代化而奋斗。《工程控制论》这一新版的作者们，正是在这一时期锻炼成长起来的中国青年控制理论科学家们。他们，尤其是宋健同志，带头组织并亲自写作定稿，完成了工作量的绝大部分，是新版的创造者。有他们这一代人，使我更感到实现四个现代化有了保障。对这一新版，我是没有做什么工作的，但为了表达对他们的敬意，同时也算是对我国20多年来伟大变革的纪念，纪念我们这一段共同的经历，我要为宋健等同志创造的新版写一篇序。

序的总题目，就是如何加速实现党中央号召，全国人民所向往的农业、工业、国防和科学技术现代化。实现四个现代化就必须发展生产力；而发展生产力的一个重要方面就是推进技术革命。所以，我就从技术革命讲起，最后说到本书的题目：控制论。

\section*{（一）}
讲技术革命，首先要提一下其他几个有关的词汇。

20世纪现代科学技术的伟大成就，正在对生产以及整个社会产生着巨大的冲击。有人常常用“新的工业革命”、“第二次工业革命”、“第三次工业革命”、“科学技术革命”等等词句来表达现代科学技术伟大成就的社会意义。但是，我们在使用这些词句时，不应忽视这些词汇的背景。

这就有必要回溯到20世纪40年代末，对这些提法的来历作一番考察。

控制论的奠基人N. 维纳在1947年10月这样说过：“如果我说，第一次工业革命是革‘阴暗的魔鬼的磨房’的命，是人手由于和机器竞争而贬值；……那么现在的工业革命便在于人脑的贬值，至少人脑所起的较简单的较具有常规性质的判断作用将要贬值。”\footnote{N. 维纳：《控制论》，科学出版社，1962年，第27页。}因此，维纳是第一个把控制论引起的自动化同“第二次工业革命”联系起来的人。此后，J. D. 贝尔纳在1954年也提出自动化是一次“新的工业革命”，他说：“我们有理由提到一次新的工业革命，因为我们引用了电子装置所能提供的控制因素、判断因素和精密因素，还有进行工业操作的速度大大增加了。巨型的自动化生产线，甚至完全自动化的工厂都有了……”\footnote{J. D. 贝尔纳：《历史上的科学》，科学出版社，第471页。}贝尔纳同时提出了“科学技术革命”这个名词，他说：“20世纪新的革命性特征不可能局限于科学，它甚至于更寄托在下列事实，就是只有在今天科学才做到控制工业和农业。这场革命或许可以更公允地叫作第一次科学—技术革命”\footnote{J. D. 贝尔纳：《历史上的科学》，科学出版社，第752页。}。在维纳和贝尔纳之后，资本主义国家的学者日渐增多地采用这两个词，而尤以“第二次工业革命”这个词更为流行。

苏联学术界在1955年以前的一段时间内，曾经把“第二次工业革命”和“科学技术革命”作为美化资本主义的概念而加以拒绝；60年代初，态度发生转变，开始接受这两个概念；到了70年代，“科学技术革命”已经成为今天苏联学术界普遍接受的概念了\footnote{费多谢耶夫：“科学技术革命的社会意义”，苏联《哲学问题》杂志，1974年第7期。}；虽然对“第二次工业革命”这个概念还有争论，但把它作为一个新概念接受下来也已成为事实。

当然，概念上的紊乱也是存在的：诸如一面讲“自动化是新的工业革命”、“计算机在工业上的应用正在引起第二次工业革命”、“第二次工业革命在本世纪早期始于美国，它指在例行的重复性的工作中，用自动控制和逻辑装置代替人的智力和神经系统”，然后又说什么“空间时代是工业革命的第三阶段。”一面讲“现代只有在苏联才发生新的工业革命”，另外又讲“不论在社会主义国家还是在发达的资本主义国家都正在发生新的工业革命即第二次工业革命”等等。在“科学技术革命”问题上的说法也很类似，诸如一方面讲“新的工业革命即科学技术革命”、又讲“科学技术革命作为一个过程，按其内容和本质是不同于工业革命的”，还说“科学技术革命是第二次工业革命的先驱”、“科学技术革命即管理工艺过程的革命”、“科学技术革命是由于科学起着优先作用而实现的现代社会生产力的根本变革”等等。其实科学技术革命这个词就容易和概念上完全不同的科学革命混淆，科学革命是指人类认识客观世界的重大飞跃，在自然科学领域里的科学革命已经由库恩\footnote{T. S. 库恩：《科学革命的结构》，上海科学技术出版社。}作了详细的阐述。所以科学革命只是认识客观世界，还不是改造客观世界，它们有联系、但又是不相同的。

苏联学术界对待“第二次工业革命”和“科学技术革命”这两个概念的态度，为什么有一个曲折的过程？这也是一个值得思考的问题。1972年，苏联《哲学问题》杂志在第12期的社论中宣称，“科学技术革命”使“生产的相互关系”、“社会的状态”和“社会的结构”等等“发生了根本的变化”，“现代世界发生的深刻变化”迫使人们对马列主义基本原理做“这样或那样的修正”\footnote{苏联《哲学问题》杂志：“今日之历史唯物主义：问题与任务”，1972年第10期社论。}。从这个观点，人们不难看出，这些人提出科学技术革命的目的是要修正马克思主义基本原理。

科学的社会科学，应该把它所有的概念同马克思主义的基础协调起来，并且实现精确化。对“第二次工业革命”、“科学技术革命”这些流行概念给以必要推敲和订正，这不仅是科学的社会科学工作者的任务，而且也是自然科学技术工作者的任务。为此就有必要回到“产业革命”或所谓“第一次工业革命”这个问题上来。

\section*{（二）}
最先提出“产业革命”概念的是恩格斯。继恩格斯之后，有法国人著作中的“产业革命”概念，也有英国资产阶级经济历史学家托因比的“产业革命”概念\footnote{A. Toynbee, \textit{Lectures on The Industrial Revolution in England}. London, (1884).}。必须说，只有马克思主义的“产业革命”概念才是真正科学的概念。恩格斯在1845年出版的《英国工人阶级状况》一书中，关于“产业革命”的论述，是科学的社会科学对“产业革命”概念的最早论述。恩格斯说：“英国工人阶级的历史是从18世纪后半期，从蒸汽机和棉花加工机的发明开始的。大家知道，这些发明推动了产业革命，产业革命同时又引起了市民社会中的全面变革，而它的世界历史意义只是在现在才开始被认识清楚。”“产业革命对英国的意义，就像政治革命对于法国，哲学革命对于德国一样。……但这个产业革命的最重要的产物是英国无产阶级”\footnote{恩格斯：《马克思恩格斯全集》，第二卷，人民出版社，1957年，第281、296页。}。

什么是“技术革命”呢？首先给予“技术革命”概念以精确化定义的是毛主席。毛主席在50年代就使用过技术革命这个词，它往往是和技术革新并列的。但毛主席没有停留在这样一般的认识上，后来他进一步发展和总结了历史上生产力发展的规律，阐明了技术革命这一概念，指出：“对每一具体技术改革说来，称为技术革新就可以了，不必再说技术革命。技术革命指历史上重大技术改革，例如用蒸汽机代替手工，后来又发明电力，现在又发明原子能之类。”毛主席举出了三个技术革命的例子，其中两个是历史上的，一个是现代的。把它们作为技术革命的典型加以研究，会给我们什么启发？毛主席这段话的历史意义和现实意义是什么呢？

蒸汽机技术革命同18世纪工业革命既有联系，又有区别。在工场手工业时期，1688年英国人托马斯·萨弗里发明了利用蒸汽冷凝产生的真空和蒸汽压力工作的抽水用蒸汽泵；两年之后法国人巴本证实了德国人莱伯尼兹提出的蒸汽可在汽缸中推动活塞的原理；1712年英国人托马斯·纽柯门做成用蒸汽和空气压力工作的一种蒸汽泵，用于矿井抽水。但是，这些蒸汽机并没有引起工业革命，相反的，正是由于创造了工具机，才使蒸汽机的革命成为必要\footnote{马克思：《马克思恩格斯全集》，第二十三卷，人民出版社，1972年，第412页。}。1764年出现珍妮纺纱机，1767年出现水力纺纱机，1785年出现骡机，这一系列工具机的发明促使瓦特实现了蒸汽机的革命。1764年他在格拉斯哥大学修理纽柯门机器的模型时产生了他的伟大发明，1769年他获得第一种蒸汽机专利，1784年获得第二种蒸汽机专利，1785年蒸汽机开始用来发动纺纱机，1786年建成博尔顿·瓦特蒸汽机工厂。

瓦特的蒸汽机是大工业普遍应用的第一个动力机，它取代了在生产过程中作为动力提供者的人。一台蒸汽机推动许多台工具机，形成有组织的机器体系，这就是工厂制度的诞生。从1786年到1800年，瓦特的工厂共生产了500多台蒸汽机，大大加速了工业革命的步伐，“工场手工业时代的迟缓的发展进程变成了生产中的真正狂飙时期”\footnote{恩格斯：《反杜林论》，人民出版社，1970年，第258页。}，蒸汽机成为大工业迅速发展的推动力，“推动力一旦产生，它就扩展到工业活动的一切部门里去……当工业中机械能的巨大意义在实践上得到证明以后，人们便用一切办法来全面地利用这种能量”\footnote{恩格斯：《马克思恩格斯全集》，第二卷，人民出版社，1957年，第291页。}。所以是蒸汽机技术革命导致了工业革命或产业革命。

我们再看毛主席举的技术革命的第二个例子：电力的发明和应用。1831年法拉第对电磁定律的发现，为电力的发明奠定了基础。1878年，爱迪生发明能在商业上普遍应用的双极发电机，并提出由一个公共供电系统向用户供电的计划；次年，爱迪生制成白炽灯；再下一年，爱迪生的电灯首先展示在“哥伦比亚”号轮船上。适应社会对这种前所未有的干净、明亮的照明工具的需要，很快出现了一个完全新型的工业——电力工业。1882年爱迪生的发电厂和供电系统在纽约运转；同一年在慕尼黑电气展览会上，法国物理学家马塞尔·德普勒展出了他在米斯巴赫至慕尼黑之间架设的第一条实验性输电线路，从此开始了交流电远距离传输技术的大发展。电力的发明从照明开始，但由于它解决了动力的分配、传输和转换问题，所以很快在大工业中得到普遍应用。马克思在他逝世前夕，曾以极为喜悦的心情密切注视着电力的发明。在1883年，恩格斯针对电力的发明说：“这实际上是一次巨大的革命。蒸汽机教我们把热变成机械运动，而电的利用将为我们开辟一条道路，使一切形式的能——热、机械运动、电、磁、光——互相转化，并在工业上加以利用。循环完成了。德普勒的最新发现，在于能够把高压电流在能量损失较小的情况下通过普通电线输送到迄今连想也不敢想的远距离，并在那一端加以利用——这件事还只是处于萌芽状态——这一发现使工业几乎彻底摆脱地方条件所规定的一切界限，并且使极遥远的水力的利用成为可能，如果在最初它只是对城市有利，那么到最后它终将成为消除城乡对立的最强有力的杠杆。但是非常明显的是，生产力将因此得到极大的发展，以至于资产阶级对生产力的管理愈来愈不能胜任。”\footnote{恩格斯：《马克思恩格斯选集》，第四卷，人民出版社，1972年，第436页。}实践证实了恩格斯的科学预见，“电力工业是最能代表最新的技术成就和19世纪末、20世纪初的资本主义的一个工业部门”\footnote{列宁：《列宁选集》，第二卷，人民出版社，1972年，第788页。}，而且直到今天也仍然是如此。是电力技术革命推进了资本主义转入垄断阶段，出现了资本帝国主义，即帝国主义。

蒸汽机和电力这两个历史上的技术革命例子，使我们把科学技术的发展作为一种社会过程、社会现象来研究，从它的发展规律，能够找到一条线索：生产力的发展史是以技术革命划分阶段的。这是毛主席关于技术革命的重要论述对我们的启示。

\section*{（三）}
生产力始终处在发展过程中，而这种发展过程又首先是从生产技术的改革开始的。生产力的发展水平取决于生产技术的高低。生产力的发展同一切事物一样，总是采取两种状态，即相对稳定的发展状态和飞跃变动的发展状态，换句话说，即生产力的发展呈现一种阶段性。“由粗笨的石器过渡到弓箭，与此相联系，从狩猎生活过渡到驯养动物和原始畜牧；由石器过渡到金属工具（铁斧、铁铧犁等等），与此相适应，过渡到种植植物和耕作业；加工材料的金属工具进一步改良，过渡到冶铁风箱，过渡到陶器生产，与此相适应，手工业得到发展，手工业脱离农业，独立手工业生产以及后来的工场手工业生产得到发展；从手工业生产工具过渡到机器，手工业-工场手工业生产转变为机器工业；进而过渡到机器制造，出现现代大机器工业，——这就是人类史上社会生产力发展的一个大致的、远不完备的情景。”\footnote{斯大林：《辩证唯物主义和历史唯物主义》，《斯大林文选，1934—1952》，上册，人民出版社，1962年，第199页。}这是对生产力发展阶段性的最形象描述。在这一幅生产力发展的大致情景中，每当生产力出现一次飞跃变动，就意味着某一技术革命被引进到了社会生产之中。革命就是量变到质变的飞跃，每一技术革命的本身就是经过一个时期实践经验的累积，有时还要经历很长的孕育时期，然后才显示出来。它一出现又立即影响了整个社会生产，引起生产力的飞跃发展。人类社会生产力和整个社会的发展就是这样波浪式地前进。技术革命是那些引起生产力飞跃发展的技术变革，不是生产力持续发展的一般技术改革或技术革新。

技术、技术革命属于劳动过程或生产过程\footnote{马克思：《马克思恩格斯全集》，第二十四卷，人民出版社，1972年，第44、123页。}，“生产过程可能扩大的比例不是任意规定的，而是技术上规定的。”\footnote{马克思：《马克思恩格斯全集》，第二十四卷，人民出版社，1972年，第91页。}技术革命乃是生产力发生飞跃变化的技术根源，而生产力的飞跃发展又必然推动社会历史的阶段性变化。是蒸汽机技术革命带来了产业革命这一生产力的飞跃变化，推进了自由资本主义的兴起；而电力技术革命却加速了资本主义的历史进程，促使它进入垄断资本主义。这就是毛主席提出技术革命这一科学概念的伟大而深远的涵义。同时，我们也看到，“技术革命”一词比前几节中介绍的其他几个词汇更精确，更有利于讨论研究问题。

为了极大地提高我国社会生产力，我们应该深入研究当前出现的几项技术革命的涵义，探索正在酝酿、即将出现的技术革命，能动地推进技术革命，加速我国四个现代化的建设。

\section*{（四）}
我们先讨论核能技术革命。

核能技术是本世纪初物理科学的伟大产物。1909年，爱因斯坦发现了质能等效性原理，预示了原子核反应所释放的能量比化学反应释放的能量大几百万倍的可能性。此后，科学家为敲开核能宝库的大门进行了不懈的努力。1932年恰德威克发现中子，找到了分裂原子核的钥匙；1938年末，哈恩和斯特拉斯曼用中子轰击铀，发现了铀原子核的可裂变性；次年1月27日，在美国华盛顿举行的物理学家会议上，波尔和费米介绍了上述发现的重大意义，费米首先提出了链式反应的理论；1942年12月2日，费米在芝加哥大学建成第一个原子反应堆，首次用实验证明：在可裂变的铀核中能够产生自维持的链式反应，从而迎来了核能技术的黎明。

核能是一种十分集中的新能源。1公斤铀所含的裂变能量约相当于2000吨煤。全世界有丰富的铀储量，在煤、石油、天然气日益枯竭的情况下，原子核的裂变能是一个有广阔前途的新能源；一座100万千瓦的核电站正常运行一年，节省的矿物燃料相当于145万吨石油或236万吨煤或16.5亿立方米天然气；据1976年国外统计资料，核电站的每千瓦小时电总平均费用已低于烧煤和石油的火力发电。正如蒸汽机出现时的情形一样，当核能的巨大意义在实践中得到证明之后，人们就会全力以赴把这种现代生产力发展的巨大推动力扩展开来。自1959年出现第一座商用核电站以来，新兴的核能电力工业正在迅速发展，截至1978年6月30日止，全世界已建成运行的电功率在3万千瓦以上的核电站已达207座，总电功率约达1.08亿千瓦；全世界正在建设的核电站有219座，总电功率达1.96亿千瓦；正在计划建设的核电站有123座，总电功率达1.24亿千瓦。预计到公元2000年，全世界核电站的装机容量将达13～16亿千瓦，届时将占全世界总发电量的45\%。

早在发现核裂变前，科学家就了解到，包括太阳在内的恒星其持续发射的巨大能量来自轻元素的核聚变。但这种核聚变反应是在一个极高的温度和压力环境中维持的，人工创造这样一个环境现在还做不到。比较容易一点的是氘的聚变，而地球海洋里就有极大量的氘；一升海水中就可以提取约33毫克的氘，这一点氘的聚变能量就等于300升的汽油！但就是核聚变也不是轻而易举的，在1945年6月15日首次裂变原子弹爆炸实现后，到1952年才实现了第一次聚变“氢弹”爆炸。现在人们正向可控制聚变反应——建设聚变反应核电站的目标前进，有可能在本世纪末实现。这样仅海水中的氘所含的能量就够人类用了。

\section*{（五）}
对现代生产产生深远影响的第二项技术革命是电子数字计算机。

蒸汽机和电力实现了生产过程的机械化，而监督与调整生产过程的工作仍需人工来完成。工人要不断照料机器的动作，用眼、耳和神经系统来直接获取生产过程的信息，然后由大脑对这些信息进行处理，作出要不要改变机器运行状况的决定，并通过手对机器的直接调整来执行这一决定。20世纪初以来，产生了能对各种物理量进行精确测量的感受器件，也产生了各种执行机构。获取机器生产状况的信息的工作，就由感受器件取代了人的器官；控制决定的执行，由执行机构取代了手对机器的直接调整。但是，控制决定还得由人直接作出，整个生产过程还需人的直接参与。这样一种状况影响着生产率进一步发展。对一些日益精密化、快速化的现代工业过程（如化学工程过程），人工控制已完全不能胜任，因为在这种情况下人的思维在速度、可靠性和耐力方面都显得不够。50年代出现了模拟式自动控制设备，在一些不太复杂的生产过程中实现了自动控制。但是，这种设备一般不能用于复杂的现代化工业过程，不能进行数据处理，也不能用于整个工厂或车间的全盘自动化。电子计算机的出现并应用于工业生产，才使自动控制技术产生了革命。第一，电子数字计算机具有计算精确的特点，和数字化感受器件、数字化执行机构结合，能够实现工业生产过程的精密控制；第二，电子数字计算机具有很大的计算能力，可以根据生产过程运行状况的改变而自动改变调节参数；可以计算出生产过程的发展趋势，以便决定应当预先调整那些操作条件。所以计算机能够对复杂的工业生产过程实现自动控制；第三，计算机不仅能对生产过程进行最优控制，而且能对包括感受器件、执行机构和计算机本身在内的全部生产设备进行监督控制。所以计算机能够实现整个企业和企业体系生产过程的全盘自动化。

关于过程的信息，是调节与控制这个过程的手段。人和人需要交换信息，人和机器也需要交换信息，任何社会实践过程都需要处理信息。人处理信息的能力，直接影响着他调节与控制事物的能力。电子计算机作为最具普遍意义的信息自动化处理设备，除了用于生产过程的数字自动化控制外，还广泛用于军事技术、科学研究、天气预报、交通运输、组织管理、信息管理、财政贸易和日常生活等领域，并成为现代化社会一种最富有代表性的装备。据1976年年底的统计数字，每百万就业人口（不包括农业）所拥有的通用电子数字计算机，美国是1 800余台。日本、联邦德国是800余台，这个数字还在迅速增长中。

\section*{（六）}
对现代生产和现代科学技术的面貌产生深远影响的第三项技术革命是航天技术。航天技术，是把航天应用于生产、科学技术和军事的一大类新技术的总称，是20世纪50年代诞生的重大技术成就。航天技术短短20余年的发展历史，不仅表现在军事上的重要性，而且显示出了它在社会生产和科学技术范围内的巨大应用潜力。

航天技术首先把作为社会生产过程一般条件的通信手段提高到了一个全新的发展水平，实现了一种理想的天上中继站——通信卫星。利用卫星通信，不需要敷设电缆或微波接力站，极少受大气干扰，作用范围广，可靠性高，而且通信容量大。一颗通信卫星的通信能力与100条越洋海底电缆相当。利用卫星通信，实现了电视对广大用户的直接广播；利用卫星通信，可以把大范围的信息处理设备沟通形成信息网络。

航天技术实现了气象观测方式的革命。气象卫星能够在全球范围内对海洋、大陆和大气层进行观测；能够昼夜提供全球性的云图照片；能够对关键性的气象参数的垂直剖面图进行精确探测；能够连续监视大片地区的天气现象，并对研究台风一类灾害性天气现象有很大的作用。航天技术还能用于监视地壳的活动现象，并对地震和火山活动的预报作出贡献。

“运输业是一个物质生产的领域”，海运、空运的导航技术与社会生产的发展紧密相关。航天技术提供了一种理想的天上无线电导航台——导航卫星，从天上直接给飞机、船舶、潜艇传送导航信号，大大提高了导航系统的经济性、可靠性和精确性。卫星导航技术的最新发展，将可以提供全球性的、连续性的、高精度的导航业务，定位误差不超过10米，测速精度为每秒3厘米，比地面无线电导航提高近100倍！

航天技术提供了一种经济、有效的自然资源大面积普查手段。地球资源卫星可以用于土壤资源的调查、规划和开发，农作物长势和病害预报，矿物资源普查，水文勘测，林业、牧业资源管理，海洋资源调查，等等。

航天技术还开辟了“天上生产”的远景。例如，在赤道同步卫星轨道上，太阳产生的能量密度约为每分钟每平方厘米二卡，而且不受地球昼夜和天气变化的影响，我们可以设想在未来利用这种环境在天上建设大型太阳能电站持续发电，然后通过大功率微波器件转换成微波能量，定向发射回地面接收站，再转换成工业和民用所需的电力。

由于航天技术的最新发展，在20世纪80年代将出现一种先进的可往返使用的航天运载工具——航天飞机。航天飞机将取代先前一次使用的卫星运载火箭；将能够对在轨道上运行的通信卫星、导航卫星、地球资源卫星、气象卫星和科学卫星进行维修服务；能把已在轨道上完成了任务的有效载荷取回地面，以便修复使用或供改进技术用；将能为航天技术提供经济的“天上实验室”；将能使利用天上无重力环境进行“天上生产”成为现实。航天飞机的发展将把航天技术革命进一步推向深入。

航天技术对生产和科学技术的发展将继续产生深远的影响。从根本上说，这是由于航天具有极其深刻的认识论意义。任何知识的来源，在于人的肉体感官对客观外界的感觉。因此，任何技术的发展都与人类眼界扩大的程度相关。航天技术提供了一个极其优越的位置，从天上来发展我们对地球、大气层和整个自然界的认识，使人的眼界有一个飞跃的扩大。在航天技术出现之前，人局限在地球上，眼界很小，对范围极其辽阔的陆地、海洋、大气层进行一番系统的考察，所需要的时间是十分长的；对范围很大的区域性、洲际性甚至全球性的自然现象，根本无法直接观察；对环境条件恶劣地区的自然现象，难以深入考察；对一些迅速变化的自然现象，人也缺乏连续观察的能力。航天技术从根本上改变了这种状况。应用目前已经成熟的技术，从数百公里高的卫星轨道上对地球拍照，一张用于地质普查的卫星照片可以覆盖地面3.4万平方公里，为普通航空观测照片的340倍！应用离地面3.5万多公里的赤道同步卫星，可以连续“俯视”大约半个地球表面。航天技术给我们提供了多种多样的天上观察站，以发展我们对自然界的认识：利用极地轨道卫星，可以在十多天内普查全球一次；利用赤道同步轨道卫星，可以连续不断监视地面自然现象；利用太阳同步轨道卫星，可以在太阳光照基本一致的条件下对自然特征进行对比研究。航天遥感技术还扩展了人对地表、洋面和大气层辐射的电磁波谱的识别范围，使一些表现在可见光区域以外的自然现象成为可以观察的。航天技术极大地延伸了人的眼力。以天文观察为例，最近发现的发射X-光和$\gamma$射线的星源和与其相关的一系列所谓高能天文学现象，没有天文卫星这个工具是不可设想的。又如从地球用光学望远镜观察火星表面只能辨认出尺度大于300公里的特征；而飞往火星的航天探测器，能在几千公里的近距离拍摄火星照片并传回地球，使分辨能力一下子提高了100倍！环绕火星的航天探测器进一步把这一能力提高到1000倍以上；而在火星表面软着陆的航天探测器则能对其表面进行直接探测，并将结果传回地球。各种各样的行星探测器使人的眼力一下子延伸了数千万甚至成亿公里！

以前我们是局限于地球表面来搞科学实验的，但就在这样的条件下，我们创造了如此丰富的科学技术，如此丰富的知识宝库。今后我们可以跳出地球表面，进入太阳系的空间，我们对宇宙的认识必然会有一个飞跃！

\section*{（七）}
除了以上所说的三项当代技术革命，核能技术革命、电子计算机技术革命和航天技术革命之外，我们还看到现代科学技术的重大突破正酝酿着另外几项技术革命。例如激光技术的发展将会导致新的技术革命，开创光子学、光子技术和光子工业\footnote{钱学森：“光子学、光子技术、光子工业”，《激光》，1979年第1期。}。又如遗传工程的发展也将会导致新的技术革命，开创按人的计划，创造新的生物种属，而不光是靠老天爷培育生物种属。还可能有其他技术革命。五、六项技术革命同时并进，百花齐放，万紫千红，是人类历史上从未有过的局面！

但所有这些科学技术的发展，所有这些技术革命都直接与控制论联在一起。控制论的发生可以追溯到电力驱动技术，即电力技术革命；而控制论的成长则同当代几项技术革命分不开的。可以预言，控制论的进一步发展也必将同我们以上论述的技术革命的进一步发展紧密相配合。让我们看一看几十年来的历史。

1944年那一台名叫 MARK-1 的大型继电器式计算机，1945年宾夕法尼亚大学那台采用电子管代替继电器的 ENIAC 电子计算机，都出现在控制论完全形成之前。但是，用替续的开关装置和用二进制作为电子计算机设计的最合适基础，完全是受惠于从1942年前后开始的控制论思想的发展：人的神经系统在做计算工作时，作为计算元件的神经元或神经细胞，实质上可以看作只具有两种动作状态的替续器。工程控制论出现以后，已日益深刻地被应用于指导电子计算机的设计。例如，能够记住主题并把以后接受的信息同这个主题联系起来的智能终端，能够识别语言波形、完全按照声音来操作的计算机，能够直接把图像转变为数字信息存储、处理的计算机，以及具有一定自学习、自组织功能的以电子计算机为心脏的机器智能等等，都是按照控制论原理来革新电子计算机体系结构的一些新发展。工程控制论正在推动电子计算机技术革命的深入。这样一个现实已经来到了人类的面前：由电子计算机和机器智能装备起来的人，已经成为更有作为，更高超的人！

工程控制论在其成形的时候，就把设计稳定与制导系统这类工程实践作为主要研究对象。虽然，作为现代火箭技术和航天技术萌芽的 V-2 火箭在控制论诞生之前好几年就出现了，但是，同应用工程控制论所实现的高精度、高可靠性的制导技术比较起来，V-2 的机电式制导系统实在是太原始了。法西斯德国向伦敦发射了2 000枚这种射程300公里的火箭，只有1 230枚落入市区，这其中又只有半数落在距目标中心13公里的范围之内。而现代制导技术可以达到这样的成就：射程1万公里的洲际导弹弹头落点圆公算偏差在30米以内；“海盗号”航天飞行器在远距地球7 000万公里之遥的火星实现了准确的软着陆。各种人造地球卫星、行星探测器、运载火箭以及航天飞机，都是高度自动化的机器。航天遥测、航天遥控、航天遥感，还有航天测控信息的远距离传递，都是工程控制论在航天技术革命发展过程中建立的里程碑。

高精度、高可靠性自动调节、自动控制和自动监测系统，对核能技术的发展极具重要性。在核电站发展的早期，一般采用常规的机电自动控制技术和仪表。1963年，在核电站调节、控制与监测工作中首次引用了电子计算机控制，并获得了很大成功；到60年代末期，电子计算机控制已在核电站上广泛应用。全面采用电子计算机监视和控制，是当前核电站技术发展的显著特征。现代电力网建设，要求核电站在运行过程中能随电网负荷的变动而自动调整功率输出，只有应用多变量最优控制以及能预测控制变量的前馈控制等现代控制理论，才能实现这个目标。

控制论的对象是系统。所谓系统，是由相互制约的各个部分组织成的具有一定功能的整体。一个蒸汽机自动调节器是一个系统，一部自动机器是一个系统，一个生物体是一个系统，一条生产线是个系统，一个企业是个系统，一个企业体系是个系统，一项科学技术工程是个系统，一个电力调节网是个系统，一个铁路调度网是个系统；还有，一个经济协作区是个系统，一个社会组织也是一个系统。有小系统，有大系统，也有把一个国家作为对象的巨系统；有工程的系统，有生物体的系统，也有既非工程的，也非生物的系统。为了实现系统自身的稳定和功能，系统需要取得、使用、保持和传递能量、材料和信息，也需要对系统的各个构成部分进行组织。生物系统的组织是一种自组织，能够根据环境的某些变化来重新组织自己的运动的工程系统是自动控制系统。

在工程系统的实践经验基础上，20世纪60年代兴起一类新的工程技术，即系统工程\footnote{钱学森、许国志、王寿云：“组织管理的技术——系统工程”，《文汇报》，1978年9月27日。}，系统工程已从工程的系统推广应用到了非工程的系统，从工程系统工程发展到了经济系统工程和社会系统工程(简称社会工程)\footnote{钱学森、乌家培：“组织管理社会主义建设的技术——社会工程”，《经济管理》，1979年第1期。}。系统工程是各类系统的组织和管理技术。各类系统工程的共同理论基础是运筹学。但控制论研究系统各个构成部分如何进行组织，以便实现系统的稳定和有目的的行动，所以系统工程又与控制论有关。这就扩大了控制论概念的影响。

另一方面也还有这样的情况：由于机械自动调节与控制技术的发展，20世纪40年代末正式形成了控制论科学。控制论原理已成功地应用于工程系统、生物系统和高级神经系统，50年代诞生了工程控制论和生物控制论。60年代，现代控制论发展形成的大系统理论，已把控制论的方法推广到了既非工程又非生物的系统——经济系统，从而正在出现一个新的控制论分支——经济控制论。面临这样一种发展形势，人们自然要问，控制论方法能否对比大系统更大的巨系统即社会系统发挥效用？

维纳在1948年曾经说过，那种认为控制论的新思想会发生某种社会效用的想法是“虚伪的希望”，“把自然科学中的方法推广到人类学、社会学、经济学方面去，希望能在社会领域里取得同样程度的胜利”，这是一种“过分的乐观”\footnote{N. 维纳：《控制论》，科学出版社，1962年，第162、163页。}。控制论的现代发展证明维纳1948年的观点是过于保守的。把一些工程技术方法推广应用到社会领域也不是“过分的乐观”，而是现实。运筹学已用于经济科学，并将应用于更大的社会领域。

恩格斯曾经预言，在社会主义条件下，“社会生产内部的无政府状态将为有计划的自觉的组织所代替”\footnote{恩格斯：《反杜林论》，人民出版社，1970年，第279页。}。充分利用社会主义经济规律的调节作用，能够组织自觉运转的经济系统，这样的系统实质上也是一种自动系统；充分利用社会主义建设的客观法则和统计规律的调节作用，如恩格斯所预言，可以实现社会生产的“有计划的自觉的组织”，实质上这就是一种巨型的系统，所以，控制论所研究的系统的运动形式，在高级形态的系统——社会系统中，也是存在的。因此，没有理由认为控制论的社会应用是一种“虚伪的希望”。这是一种已经看得见曙光的真实的希望。在社会主义条件下，一门新的科学终将诞生，这就是社会控制论。这样一门科学不会在资本主义制度下出现，因为，“资产阶级社会的症结正是在于，对生产自始就不存在有意识的社会调节”\footnote{恩格斯：《马克思恩格斯选集》第四卷，人民出版社，1972年，第369页。}。

\section*{（八）}
作为技术科学的控制论，对工程技术、生物和生命现象的研究和经济科学，以及对社会研究都有深刻的意义，比起相对论和量子论对社会的作用有过之无不及。我们可以毫不含糊地说从科学理论的角度来看，20世纪上半叶的三大伟绩是相对论、量子论和控制论\footnote{童天湘：“控制论的发展和应用”，《哲学研究》，1979年第2期。}，也许可以称它们为三项科学革命，是人类认识客观世界的三大飞跃。但我们比较这三大理论，也看到它们，特别是前两者与后者的区别。相对论是处理宏观物质运动的基础理论，量子力学是处理微观物质运动的基础理论。它们有一个共同点，都是研究物质运动的；还有一个共同点，都是基础理论，即人们实践的基本总结，物质运动不管什么形式，都是以此为依据的。控制论则不然，它的研究对象似乎不是物质运动，而且好像还没有深入到可以称为基础理论。这就发人深思了。相对论和量子力学的典型可以引导我们设想控制论的进一步发展的方向。

为什么说控制论似乎还不够深入呢？从控制论上述的形成和发展来看，它是原始于技术的，即从解决生产实践问题开始的。工程控制论首先建立，是控制工程系统的技术的总结，即从工程技术提炼到工程技术的理论，即技术科学。有了这样一门技术科学——工程控制论，就如前面讲到的，我们又发现生物生命现象中的一些问题也可以用同样的观点来考察，从而建立了生物控制论。再进而发展到经济控制论以及社会控制论。现在我们如果把这四门技术科学加在一起称为控制论，这样形成的所谓控制论还是一个混合物，没有脱离其本来技术科学的面目；特性的内容多些，普遍存在的共性内容不够突出。能不能更集中研究“控制”的共性问题，从而把控制论提高到真正的一门基础科学呢？能不能把工程控制论、生物控制论、经济控制论、社会控制论等等，作为是由这门基础科学理论控制论派生出来的技术科学呢？

理论控制论的对象是不是物质的运动？因为世界是由运动着的物质构成的，控制论的对象自然还是客观世界，所以控制论的研究对象最终还得联系到物质，只不过不是物质运动本身而是代表物质运动的事物因素之间的关系。有些关系是直接的；有些关系不直接，要通过信息通道，表现为信息。此外，为了控制，即使受控对象按我们的预定要求行事，我们还加入若干控制量和控制量与事物因素以及信息之间的关系。事物因素、信息和控制量形成一个相互关联体系，表现为可以用数学表达的一系列关系。我们要注意关联必须以数学形式定下来，也就是要定量，不然就没有控制论。理论控制论的任务就是根据这些定量的关系预见整个系统的行为。有些问题在控制论中是有决定性意义的：如系统的能控性问题和能观测性问题的普遍理论。

如果这就是我们要建立的基础科学理论控制论，那我们可以从这一新版《工程控制论》看到，我们达到的离我们的目标还有一定距离，深度还很差。要真正建立这门基础科学，还有待于今后控制论专业工作者们的努力。为了实现我国的社会主义现代化，为了促进当前的和即将到来的各项技术革命，我认为这一努力是很有意义的。

在写这篇序的过程中，王寿云同志帮助我查阅并整理了很多资料，付出了辛勤的劳动，我在此对他表示感谢。

\begin{flushright}
写于1978年12月24日\\
修改于1979年11月29日
\end{flushright}